%% CV — Ho Won KIM
%% pdfLaTeX · A4

\documentclass[10.5pt, a4paper]{article}

\usepackage[utf8]{inputenc}
\usepackage[T1]{fontenc}
\usepackage[french]{babel}
\usepackage{lmodern}
\usepackage{microtype}

\hyphenpenalty=10000
\exhyphenpenalty=10000
\emergencystretch=3em
\usepackage[
  top=1.8cm, bottom=1.8cm,
  left=2.0cm, right=2.0cm
]{geometry}
\usepackage{xcolor}
\usepackage{array}
\usepackage{tabularx}
\usepackage{booktabs}
\usepackage{parskip}
\usepackage{enumitem}
\usepackage{hyperref}
\usepackage{ragged2e}

% ── Couleurs ───────────────────────────────────────────────────────────────
\definecolor{navy}{RGB}{22, 40, 75}
\definecolor{darkgray}{RGB}{80, 80, 80}
\definecolor{lightgray}{RGB}{180, 180, 180}

\hypersetup{colorlinks=true, urlcolor=navy, linkcolor=navy}

% ── Pas d'en-tête/pied de page ────────────────────────────────────────────
\pagestyle{empty}
\setlength{\parindent}{0pt}
\setlength{\parskip}{0pt}

% ── Largeur colonne date ───────────────────────────────────────────────────
\newlength{\datecol}
\setlength{\datecol}{2.6cm}
\newlength{\contentcol}
\setlength{\contentcol}{\textwidth}
\addtolength{\contentcol}{-\datecol}
\addtolength{\contentcol}{-1.2em}  % column sep

% ── Commandes utilitaires ─────────────────────────────────────────────────

% Ligne d'entrée principale : date | contenu
\newcommand{\cvrow}[2]{%
  \noindent
  \parbox[t]{\datecol}{\raggedleft\textcolor{darkgray}{\small #1}}%
  \hspace{1.2em}%
  \parbox[t]{\contentcol}{\sloppy\hbadness=10000 #2}%
  \par\vspace{0.6em}%
}

% Section principale
\newcommand{\cvsection}[1]{%
  \vspace{10pt}%
  {\color{navy}\large\bfseries\MakeUppercase{#1}}%
  \vspace{2pt}%
  {\color{navy}\hrule height 0.8pt}%
  \vspace{6pt}%
}

% Sous-section
\newcommand{\cvsubsection}[1]{%
  \vspace{5pt}%
  {\small\bfseries\textcolor{darkgray}{\MakeUppercase{#1}}}%
  \vspace{3pt}\par%
}

% Entrée sans date (pour listes simples)
\newcommand{\cvitem}[2]{%
  \noindent
  \parbox[t]{\datecol}{\raggedleft\textcolor{darkgray}{\small #1}}%
  \hspace{1.2em}%
  \parbox[t]{\contentcol}{\sloppy\hbadness=10000 #2}%
  \par\vspace{0.6em}%
}

% ══════════════════════════════════════════════════════════════════════════
\begin{document}
% ══════════════════════════════════════════════════════════════════════════

% ── EN-TÊTE ───────────────────────────────────────────────────────────────
\begin{center}
  {\huge\bfseries\color{navy} Ho Won \textsc{Kim}}\\[10pt]
  {\large Doctorant en Sciences du Langage · Université Sorbonne Nouvelle}\\[6pt]
  {\small
    Paris, France
    \quad\textcolor{lightgray}{|}\ \
    \href{mailto:ho-won.kim@sorbonne-nouvelle.fr}{ho-won.kim@sorbonne-nouvelle.fr}
    \quad\textcolor{lightgray}{|}\ \
    \href{https://howon-kim.com}{howon-kim.com}
  }
\end{center}

\vspace{8pt}

{Doctorant au CLESTHIA (EA~7345), Université Sorbonne Nouvelle, spécialisé en analyse du
discours et en sémantique. Recherches portant sur le discours institutionnel de TotalEnergies dans le cadre de la transition énergétique, selon une approche énonciative et modale.}
\vspace{4pt}
\\
{\small\itshape\color{darkgray}
Mots clés : Analyse du discours  ·  Sémantique  ·  Pragmatique  ·  Discours numérique  ·  Discours environnemental}
\vspace{6pt}


% ── FORMATION ─────────────────────────────────────────────────────────────
\cvsection{Formation}

\cvrow{2024 -- présent}{%
  \textbf{Doctorat en Sciences du Langage}
  \textit{(depuis la 2\textsuperscript{e} année)}\\
  CLESTHIA (EA~7345) · Université Sorbonne Nouvelle\\
  {\small   · Thèse : \textit{«~Construction modale de la réalité dans le discours sur la
  transition énergétique : analyse du discours de TotalEnergies (2019--2025)~»}\\
    · Directrice de thèse : Georgeta Cislaru}}

\cvrow{2023 -- 2024}{%
  \textbf{Doctorat en Sciences du Langage}
  \textit{(1\textsuperscript{re} année)}\\
  MoDyCo (UMR~7114) · Université Paris Nanterre}

\cvrow{2020 -- 2023}{%
  \textbf{Master en Langue et Littérature Françaises}\\
  Université Nationale de Séoul\\
  {\small   · Mémoire : \textit{«~Analyse de discours climatique français :
  Dénomination et son sens lexical~»}
\\
   ·  Directeur : Seong Heon Lee}}

\cvrow{2019 -- 2020}{%
  \textbf{Étudiant en échange} · Sciences Po Paris}

\cvrow{2014 -- 2020}{%
  \textbf{Licence en Langue et Littérature Françaises}\\
  Université Nationale de Séoul\\
  {\small   · Mention \textit{Summa cum Laude}
  \\· Spécialité PPE (Philosophie--Politique--Économie)\\
  · Mémoire d'excellence}}

% ── PUBLICATIONS ──────────────────────────────────────────────────────────
\cvsection{Travaux \& Publications}

\cvsubsection{Articles dans des revues}

\cvrow{2019}{%
  Kim Ho Won, Choi Yun Son — «~Analyse critique du discours de la presse française sur
  les gilets jaunes~». \textit{Études de Langue et Littérature Françaises}
  (Revue Coréenne), 119, pp.~215--242.}

\cvsubsection{Communications dans des colloques et journées d'études}

\cvrow{2026 · 06}{%
  \textbf{«~Building a Multilingual Pipeline for the Automated Recognition of Online
  Polemics : From Marker Identification to Data Collection~»}\\
  {\small\textit{Dissemination Event «~Language Rights, Social Inclusion \& Polarised Discourses
  in Europe~» · ACVALIN · ARENAS · POLAR}
  · Université de Chypre}}

\cvrow{2026 · 06}{%
  \textbf{«~Modalité et aspectualité de la transition : les mots du développement durable dans
  les rapports de durabilité de TotalEnergies~»}\\
  {\small\textit{Journée d'étude «~Autour des mots du développement durable : Traductions, circulations, usages et effets~»}
  · Revue internationale des études du développement (RIED)}}

\cvrow{2026 · 04}{%
  \textbf{«~Constructions modales de la réalité et leurs collisions : analyse du discours de
  TotalEnergies (2019--2025)~»}\\
  {\small\textit{Journée d'études}
  · CLESTHIA (EA~7345), Université Sorbonne Nouvelle}}

\cvrow{2025 · 08}{%
  \textbf{«~Usage pragmatique des expressions honorifiques coréennes
  \textit{sugohasieossseubnida} et \textit{gosaenghasieossseubnida}~»}\\
  {\small\textit{Colloque international PHRASEOPRAG}
  · Université de Hokkaido, Sapporo, Japon}}

\cvrow{2023 · 12}{%
  \textbf{«~Durable, Vert, Renouvelable : Modification omniprésente, généralisation et
  simplification sémantique~»}\\
  {\small\textit{Journée des doctorants}
  · MoDyCo (UMR~7114), Université Paris Nanterre}}

\cvsubsection{Interventions dans des séminaires}

\cvrow{2026 · 03}{%
  \textbf{«~Construction modale de la réalité et stratégies discursives : une analyse
  linguistique des Assemblées générales de TotalEnergies (2019--2025)~»}\\
  {\small\textit{Séminaire Opération DDD}
  · CLESTHIA (EA~7345), Université Sorbonne Nouvelle}}

\cvrow{2026 · 03}{%
  \textbf{«~Fossil Fuel Industry Discourse and its Societal Circulation~»}\\
  {\small\textit{Séminaire Projet SCAIRY} · YUFE Alliance}}

\cvrow{2025 · 06}{%
  \textbf{«~Entre phraséologie et politesse : étude pragmatique de
  \textit{sugohasieossseubnida} et \textit{gosaenghasieossseubnida}
  dans le coréen contemporain~»}\\
  {\small\textit{Séminaire}
  · Société académique coréenne de philosophie à Paris}}

\cvrow{2025 · 03}{%
  \textbf{«~La nomination des énergies dans les discours corporates —
  Le cas de TotalEnergies~»}\\
  {\small\textit{Séminaire doctoral}
  · CLESTHIA (EA~7345), Université Sorbonne Nouvelle}}

% ── PROJETS ───────────────────────────────────────────────────────────────
\cvsection{Projets de Recherche}

\cvrow{2026 -- présent}{%
  \textbf{Opération «~Atlas discursif et critique du développement durable~»}
  · \textbf{Membre}\\
  {\small CLESTHIA (EA~7345), Université Sorbonne Nouvelle\\
   · Participation régulière aux séminaires et activités scientifiques de l'opération,
  incluant la présentation d'une communication de recherche.}}

\cvrow{2025 -- 2027}{%
  \textbf{Projet POLAR}
  \textit{(Polemics in Online Language and Automated Recognition)}
  · \textbf{Chargé de coordination}\\
  {\small YUFE Alliance · Université Sorbonne Nouvelle,
  Université de Chypre, Université d'Anvers\\
   · Coordination administrative et organisationnelle ; appui à la planification,
  au suivi des activités et à la mise en œuvre du programme de recherche.}}

\cvrow{2025 -- 2027}{%
  \textbf{Projet SCAIRY}
  \textit{(Sustainable Climate-Change Action via Interdisciplinary Research in YUFE)}
  · \textbf{Junior Investigator}\\
  {\small YUFE Alliance\\
   · Participation au projet interdisciplinaire autour de l'action climatique ;
  présentation sur les discours des entreprises fossiles.
  Sélectionné comme participant au BioYUFE Annual Meeting 2026.}}

\cvrow{2024 -- présent}{%
  \textbf{Opération «~Les discours du développement durable :
  circulations et tensions~»}
  · \textbf{Membre}\\
  {\small CLESTHIA (EA~7345), Université Sorbonne Nouvelle\\
   · Participation régulière aux séminaires et activités scientifiques de l'opération,
  incluant la présentation d'une communication de recherche.}}

% ── ORGANISATION ──────────────────────────────────────────────────────────
\cvsection{Organisation de manifestations scientifiques}

\cvsubsection{Colloques}

\cvrow{2026 · 12}{%
  \textbf{Colloque international
  «~The Expression of Indebtedness in Linguistics~»}\\
  {\small CEL (UR~1663) · CLESTHIA (EA~7345) · ED~661
  · European Union Next-GenerationEU\\
   · Coordination de l'organisation scientifique et logistique du colloque.}}

\cvrow{2024 · 10}{%
  \textbf{Colloque international «~ColDoc2024~»}\\
  {\small MoDyCo (UMR~7114), Université Paris Nanterre\\
   · Organisation du colloque des jeunes chercheurs et doctorants en linguistique.}}

\cvrow{2024 · 06}{%
  \textbf{Colloque international «~SigWriting 2024~»}\\
  {\small European Association for Research on Learning and Instruction (EARLI)\\
   · Organisation de l'école de recherche ; présidence de séance et modération.}}

\cvsubsection{Journées d'études}

\cvrow{2026 · 04}{%
  \textbf{Journée des doctorants} \\
  CLESTHIA (EA~7345), Université Sorbonne Nouvelle\\
  {\small  · Organisation générale ; présidence de séance, coordination et modération.}}

% ── ENSEIGNEMENT ──────────────────────────────────────────────────────────
\cvsection{Expériences d'enseignement et d'encadrement}

\cvsubsection{Chargé de cours}

\cvrow{2026--27 · S2}{%
  \textbf{Introduction à la linguistique française II} (L1, 21h · distanciel)\\
  LLFL, Université Sorbonne Nouvelle\\
  {\small  · La morphologie et la sémantique.}}

\cvrow{2026--27 · S1}{%
  \textbf{Introduction à la linguistique française I} (L1, 14h · présentiel)\\
  LLFL, Université Sorbonne Nouvelle\\
  {\small  · Introduction générale à la linguistique ; phonétique et phonologie.}}

\cvrow{2025--26 · S2}{%
  \textbf{Introduction à la linguistique française II} (L1, 21h · distanciel)\\
  LLFL, Université Sorbonne Nouvelle\\
  {\small  · La morphologie et la sémantique.}}

\cvrow{2025--26 · S1}{%
  \textbf{Introduction à la linguistique française I} (L1, 24h · présentiel)\\
  LLFL, Université Sorbonne Nouvelle\\
  {\small  · Introduction générale à la linguistique ; phonétique et phonologie.}}

\cvrow{2024--25 · S1}{%
  \textbf{Sémantique lexicale} (L3, 14h · distanciel)\\
  LLFL, Université Sorbonne Nouvelle\\
  {\small  · Créativité sémantique : métaphore et métonymie,
  déviance et lexicalisation.}}

\cvsubsection{Interventions ponctuelles invitées}

\cvrow{2025 · 07}{%
  \textbf{Alumni Talk : la vie de la recherche} (Licence, 2h)\\
  {\small SNU in Paris École d'été
  (Université Nationale de Séoul · Université Sorbonne Nouvelle)}}

\cvrow{2023 · 10}{%
  \textbf{Analyse du discours dans ses représentations sociales (Ecrifore)}
  (M1, 2h) \\ Université Paris Nanterre\\
  {\small  · Analyse critique du discours \& linguistique systémique fonctionnelle.}}

\cvsubsection{Encadrement de stagiaires}

\cvrow{2026}{%
  \textbf{Co-encadrement d'un stage de Master~2}\\
  Projet POLAR, Université Sorbonne Nouvelle\\
  {\small  · Co-encadrement avec Georgeta Cislaru · Accompagnement scientifique et méthodologique d'un étudiant de M2
  dans ses travaux de recherche en sciences du langage.}}

\cvsubsection{Assistant d'enseignement}

\cvrow{2021 -- 2022}{%
  \textbf{Initiation à l'Informatique} (L1, 60h) \\
  Collège d'enseignement libéral, Université Nationale de Séoul\\
  {\small  · TD de démonstrations Python · sessions de questions-réponses.}}

\cvrow{2018 -- 2019}{%
  \textbf{La Vie et l'Humanité} (L1, 15h) \\
  Collège des sciences humaines, Université Nationale de Séoul\\
  {\small  · Organisation et coordination du programme.}}

\cvrow{2018 -- 2019}{%
  \textbf{Ad Fontes} (Formation Continue, 60h) \\
  Collège des sciences humaines, Université Nationale de Séoul\\
  {\small  · Gestionnaire du programme de Sciences Humaines pour dirigeants en Corée.}}

% ── BOURSES ───────────────────────────────────────────────────────────────
\cvsection{Bourses \& Distinctions}

\cvitem{2023 -- 2026}{Bourse doctorale · \textit{Fondation PonyChung} (sept. 2023 -- août 2026)}
\cvitem{2021 -- 2022}{Bourse de master · \textit{Fondation Woosan} (jan. 2021 -- déc. 2022)}
\cvitem{2020 -- 2022}{Bourse d'enseignement et de recherche · Université Nationale de Séoul (4 semestres)}
\cvitem{2014 -- 2020}{Bourse de mérite · Université Nationale de Séoul (5 semestres d'excellence académique)}

% ── AUTRES ACTIVITÉS ──────────────────────────────────────────────────────
\cvsection{Autres activités}

\cvrow{2023 · 01}{%
  \textbf{Visite de recherche}
  · ATILF (UMR~7118, CNRS -- Université de Lorraine), Nancy\\
  {\small  · Activités scientifiques de lexicologie sous la direction d'Alain Polguère
  (séminaire, réunion de travail, cours).}}

\cvrow{2014 -- 2015}{%
  \textbf{Chercheur bénévole sur le harcèlement moral}\\
  {\small Avocats Coréens pour l'Intérêt Public et le Droit Humain (KLPH)\\
   · Recherche sur le harcèlement moral · traduction de codes et ordonnances juridiques
  français.}}

% ── COMPÉTENCES ───────────────────────────────────────────────────────────
\cvsection{Compétences \& Outils}

\cvitem{Langues}{Coréen \textit{(natif)} · Français \textit{(C2)} · Anglais \textit{(C2)}}
\cvitem{Corpus \& TAL}{TXM · Iramuteq · SketchEngine · Python}

% ── AUTRES FORMATIONS ─────────────────────────────────────────────────────
\cvsection{Autres formations}

\cvrow{2024 · 06}{%
  \textbf{École de printemps «~Sens de l'écologie, écologie du sens~»}\\
  {\small  PLURIELLES (UR~24142), ICAR (UMR~5191) · Lyon · Participant actif.}}

\cvrow{2023 · 08}{%
  \textbf{«~Climate Reality Leadership~»} \\
  The Climate Reality Project · Séoul}

\cvrow{2021 · 08}{%
  \textbf{«~Data Science Boot Camp~»} \\
  École supérieure de Data Science, Université Nationale de Séoul\\
  {\small   · Modules : «~Mathématiques pour la data science~»,
  «~Programmation pour la data science~».}}

\cvrow{2019 · 08}{%
  \textbf{École d'été de linguistique}
  \\ Société coréenne de linguistique · Séoul\\
  {\small   · Modules : «~Statistiques linguistiques~», «~Sémantique et pragmatique~»,
  «~Psycholinguistique~».}}

\cvrow{2019 · 09}{%
  \textbf{«~Tensorflow développeur~»}
\\   Centre de la Formation Continue, Université Nationale de Séoul}

\cvrow{2019 · 06}{%
  \textbf{«~Python et l'analyse des données~»}
\\   Centre de la Formation Continue, Université Nationale de Séoul}

\cvrow{2019 · 05}{%
  \textbf{«~Initiation au Python pour Machine Learning et Deep Learning~»}\\
  {\small   Centre pour le développement de l'enseignement et de l'apprentissage,
  Université Nationale de Séoul}}

\end{document}
